The retreat horn told her what was needed before anyone thought to send her word directly, and Melisandre understood, hearing it, that the trenches she had lit at the battle’s opening were not the trenches this night still required of her.

The kingsroad south ran through open ground for the first mile past Winterfell’s walls, ground with no natural chokepoint to slow whatever came after the retreating column once the castle itself fell — and she had known, since the war council first laid its plans across that inadequate old map, that this gap would eventually be hers to close, one way or another, when the time came.

She found Davos at the gate, directing the last wagons through with a soldier’s flat urgency that had replaced, these past hours, whatever remained of his careful courtesies. Wounded men groaned in the carts around them, and somewhere behind, deeper in the burning castle, she could hear the particular high, ragged sound of a woman calling a name that would not be answered again.

“The road,” she told him. “I can seal it. Not forever. Long enough.”

“Long enough costs you what, exactly?” He did not ask it kindly, but he asked it plainly, which she had come to understand, across their short, hard acquaintance, was the only kind of concern Davos Seaworth still had left to offer her.

“I don’t know,” she said, honestly. “I’ve never asked this much of the fire at once. I suspect I’m about to find out.”

“You said that to me once before. About Shireen. That you didn’t know, and you were about to find out.” His jaw worked, old grief surfacing plain across his weathered face even in the chaos. “I’ll not pretend I’ve forgiven you for how that particular finding-out went. But I’ll not stop you from this one either, whatever it costs you. This road needs closing more than it needs my forgiveness tonight.”

“That,” Melisandre said, “is more honesty than I had any right to expect from you, Ser Davos. I’ll take it, and be glad of it, whatever else this night still owes either of us.”

\sceneBreak

She walked the length of open ground herself, because the working required her presence rather than merely her will, red robes stark against snow already trampled by ten thousand retreating feet, and felt, raising her arms to whatever remained of the night sky, the particular vertigo of a woman reaching for more than she had ever reached for previously and finding, to her own genuine surprise, that the reach itself answered.

Fire came up out of the frozen ground in a wall higher than the one at the battle’s opening, wider, hungrier, catching along a mile of open kingsroad in a single continuous roar that lit the retreating column’s rear guard bright as noon and cast the pursuing dead, further back, into stark black silhouette against it — and Melisandre felt, holding it, her own borrowed years beginning to come due all at once, the ruby at her throat burning against her skin with a heat that had nothing to do with the flames themselves.

Her legs went first, some strength she had not questioned in three hundred years simply ceasing to answer her, and she held herself upright through will alone, arms still raised, the fire still climbing, because she understood that the moment her will faltered the wall would falter with it, and the wagons still crossing that mile of open ground would falter too, permanently, in whatever way the dead’s patience finally decided to punish the failing.

She held it. She understood, holding it, precisely what the holding was going to cost her, and found she did not resent the cost nearly so much as she’d have expected to, across a life spent believing herself entitled to considerably more time than she was now spending in a single night’s asking.

\sceneBreak

She had seen Jon Snow in her fires more times than she’d told anyone, and a shadow beside him she’d never once managed to read clearly, and had believed, across three years of watching that shadow without understanding it, that the prophecy she’d given her whole life to was finally, at long last, closing properly around the man she’d chosen to believe in.

She did not know, dying, that the shadow had never been Jon at all.

She felt the fire begin to gutter as her own strength gave out beneath it, held it a moment longer through will alone, long enough — she would never know exactly how long, only that it was enough, because the retreating column’s rearmost wagons cleared the gap before the wall finally began to fail — and felt the ruby at her throat go suddenly, completely cold, the borrowed years arriving all at once the way a debt long deferred arrives on the one day a man can no longer delay the paying of it.

She did not choose to remove the necklace. She found, in the final moment, that her hands no longer answered her at all, and understood, with something closer to relief than fear, that the choosing had already been made for her the moment she’d decided the fire was worth the asking.

\sceneBreak

Davos found her afterward, kneeling in snow gone dark with a night’s accumulated grief, an old woman’s body wearing robes cut for someone considerably younger, ancient in a way that made the red hair and unlined face she’d worn for decades seem, in hindsight, like the greater unreality rather than this withered, honest truth beneath it.

He did not weep for her. He understood, kneeling beside a woman he had never forgiven and never would, that grief and forgiveness were not, in the end, the same debt, and that he could grieve what her fire had bought the retreating column without ever once excusing what her fire had cost him at Dragonstone.

“She held it,” he said, to no one in particular, to Jon, who had come to stand beside him in the firelight’s last dying glow. “However else this night judges her. She held it long enough.”

“Long enough for what?” Jon asked, quiet.

“I don’t rightly know,” Davos said, and found, saying it, that the not-knowing was, for once, its own kind of peace. “Long enough for whatever the flames actually meant by her, I suppose, rather than whatever she believed they meant. I don’t imagine we’ll ever properly untangle the difference. I don’t imagine she managed to, either, in the end.”

They left her there, because the retreat did not permit them the luxury of carrying every body it had cost that night, and Melisandre of Asshai, priestess of R’hllor, three hundred years or considerably more beneath a face she had worn borrowed for so long she’d nearly forgotten it wasn’t hers, was given over at last to the same cold ground she had spent her whole unnaturally long life burning back.